% =====================================================================
\section{Burstiness of the load: a hypothesis that did not pass}\label{sec:burst}
Two discrepancies with the arithmetic of the hazard pointed in \emph{opposite} directions, and that looked like a lead. The observed increase of the collapse rate due to surges turned out to be tenfold \emph{stronger} than predicted; the observed fraction of cells damaged at a doubled daily range, threefold \emph{weaker}. Neither the convexity of the hazard nor the replacement of the density by the exceedance probability can produce such a thing: both corrections move the prediction in one direction. What remained was a mechanism absent from the hazard altogether --- lag of the defence: a short surge arrives at an unadapted cell, whereas a range raised for a day is partly compensated by the defence in time.

\textbf{How this was tested without a single fitted number.} To write an adaptation factor into the predictor would mean introducing a parameter derivable only from the algebra of the fold, that is, in effect fitting it. The question was therefore posed as a comparison in which everything extraneous is pinned: the strength of the slow noise sets the \emph{variance} of the load at an unchanged mean (the mean of the lognormal is compensated within the model itself), and the correlation time sets the \emph{time scale} at an unchanged mean and variance. A scan over the correlation time isolates exactly one thing: whether the defence has time to adapt.

\textbf{The rule as written is not satisfied.} The fraction of cells damaged, over ordered correlation times, gave $z=0.61$, $p=0.27$: there is no trend.

\textbf{And the defect here is mine, not the model's.} The fraction damaged --- ``was the set damaged even once in 60 days'' --- \emph{saturates}: it runs $0.599\to0.623\to0.626\to0.627\to0.614\to0.588$, that is, as a hump, and has almost no sensitivity. The collapse rate on the same blocks, meanwhile, falls monotonically from 0.0550 to 0.0328 --- by 40\,\% --- and the episode duration grows from 5.37 to 8.04 days (fig.~\ref{fig:v5}b). An estimate from the published table gives quantitative agreement as well: the share of the fall falling between the two constants of the defence equals 0.523 against the 0.424 expected under a fall uniform in the logarithm.

\begin{note}
This estimate \textbf{is not a verdict} and has not been entered in the ledger. It was computed after the data had been seen, so its threshold is fitted to them. For F24 to receive an honest verdict, the rule is rewritten in terms of the collapse rate and rerun on a \emph{different} configuration --- a different sid and a different strength of the slow noise --- with thresholds written down in advance. The verdict ``not satisfied'' of 17.10 remains in the ledger as it stands; the defect is noted beside it. This is the same discipline by which we rewrite nothing after the fact, our own misses included.
\end{note}

\textbf{The lesson is wider than the rule itself.} To the four statistical traps of version~4 a fifth was added, and it is about the choice of the \emph{measured quantity}, not of the criterion: a saturating quantity quenches a real effect, and no statistics will correct that. The mechanism acts on the rate of events --- so the rate must be measured, not the fraction of those to whom the event happened at least once.
